% Role C / kangw24. Evidence source: PR #91 full 120-job ImageNet-64 matrix.
% Keep frozen all-seed results separate from any stable-seed sensitivity view.
\subsection{Transient intervention effects and trajectory stability are distinct}
\label{sec:transient-intervention-vs-stability}

The IA--IB contrast and the stability of either trajectory answer different
questions.  At a matched seed and budget, IA--IB measures the realized-spacing
intervention under the frozen pair construction.  In PR~\#91, IA has lower
FID and KID in every one of the 30 same-seed comparisons through 6,400 kimg,
covering three seeds, five checkpoints, and NFE1/2.  Thus the repeated cells
favor IA throughout the evaluated early-to-middle horizon in each of the three
training seeds; the 30 cells are not 30 independent trajectory replications.

Trajectory stability asks whether a run continues to preserve or improve its
quality as training proceeds.  Seed 103 reaches useful quality in both arms at
7,680 kimg and then degrades in both IA and IB.  The first post-7,680 checkpoint
is especially severe for IB, but IA also deteriorates and reaches comparably
collapsed quality at later checkpoints.  Seeds 101 and 102 do not show this
pattern and continue to improve through 12,800 kimg.  The post-8,960
three-seed endpoint means are consequently highly sensitive to one unstable
paired seed rather than describing a common late-training response.

The paired endpoint contrast remains the same estimand, but its three-seed mean
is highly sensitive to whether a sampled trajectory enters this late unstable
regime.  We retain seed 103 in the frozen analysis, report its two arms
separately, and interpret the IA--IB contrast only at matched seed and budget.
Stable-seed summaries are descriptive sensitivity analyses, not replacements
for the frozen cohort.
The paired instability identifies a shared late-trajectory failure rather than
a treatment-specific cause or prevention effect.  Locating that mechanism
requires diagnostics taken on the ImageNet
trajectories before and during the transition, rather than an endpoint FID
comparison alone.
